\section{Carry history in the caller's input, not hidden kernel state}
For a chronological input of \(n+7\) samples, a second call starting at the next
output's window advances its base by \(2n\) bytes. The last seven samples of the
first call's input become the history of that next window. Coefficients retain the
same order. A caller can retain those samples in a larger stream buffer or copy them
into a new prefix; the kernel needs no persistent history object.

For a split after \(m\) outputs, the first input slice contains positions zero through
\(m+6\). The second starts at original position \(m\). Its first seven samples
are exactly those needed before the next new sample. Concatenating the outputs of
these two calls must match one complete call, including the zero-length split cases.
The checks exercise every cut from zero through sixty-four in one execution session.

Reset the per-output accumulator and per-call clipping counter at their respective
boundaries. Neither belongs in retained global state. A prior clipped call must not
change a later quiet call's count. The input and coefficient arrays remain immutable,
so this interface excludes in-place output. The exact-alias permission in the preceding
mixer chapter does not transfer to a new contract automatically.
